\section{Introdução}\label{sec:introducao}

As doenças cardiovasculares são a principal causa de morte, e o eletrocardiograma (ECG) é um meio simples, confiável e não invasivo de acompanhar a atividade do coração e de diagnosticar arritmias~\cite{zhang2021}. O ECG padrão tem 12 derivações, registradas ao mesmo tempo por eletrodos na superfície do corpo~\cite{zhang2021}. Em cada derivação, cada batimento aparece como uma sequência de ondas e segmentos, como a onda P, o complexo QRS, o segmento ST e a onda T, cujas amplitudes diferem muito entre si, com o pico R do QRS dominando o traçado~\cite{gumpfer2026}. Os diagnósticos se apoiam em partes diferentes dessa sequência: alguns na forma do QRS, de grande amplitude, como os bloqueios de ramo; outros em segmentos de baixa amplitude, próximos da linha de base, como na isquemia, que altera o segmento ST~\cite{gumpfer2026}. Como interpretar o ECG é difícil mesmo para cardiologistas experientes e o número de exames é grande, o diagnóstico automático tornou-se um tema de pesquisa ativo~\cite{zhang2021}.

Redes neurais profundas aprendem esse diagnóstico diretamente do sinal. Na base CPSC2018, com nove classes diagnósticas, uma rede residual unidimensional de 34 camadas obteve F1 médio de 0,813~\cite{zhang2021}. O desempenho, porém, não foi uniforme entre as classes: o F1 passou de 0,9 na fibrilação atrial e no bloqueio de ramo direito e foi de 0,535 na elevação do segmento ST, diferença que os autores atribuíram em parte à discordância entre médicos no diagnóstico dessa classe~\cite{zhang2021}.

Esses modelos são, em geral, considerados caixas-pretas, porque as suas decisões não podem ser rastreadas por uma pessoa, o que pesa sobretudo na prática clínica~\cite{zhang2021}. Os métodos de explicação \emph{post hoc} procuram contornar esse problema sem mudar o modelo já treinado~\cite{turbe2023}: para uma predição, atribuem a cada ponto da entrada um valor que indica quanto ele contribuiu para a saída, formando um mapa de atribuição sobre o traçado. O SHAP é um desses métodos. Ele distribui entre os atributos da entrada a diferença entre a saída do modelo e um valor de base, a saída esperada quando nenhum atributo é conhecido, de modo que a soma das atribuições reproduz essa diferença~\cite{lundberg2017}. Em redes profundas, essa distribuição é aproximada a partir dos gradientes do modelo; o GradientSHAP, por exemplo, multiplica gradientes pela diferença entre a entrada e exemplos de referência~\cite{kokhlikyan2020}. O SHAP já foi usado para explicar as predições da rede residual no CPSC2018, por paciente e por derivação~\cite{zhang2021}.

Um mapa de atribuição, contudo, não garante que descreve o modelo. Em geral, não há uma explicação correta conhecida com a qual compará-lo, e a inspeção visual é um guia fraco: mapas que não dependem do modelo nem dos dados podem ser visualmente parecidos com mapas de saliência~\cite{adebayo2018}. Por isso, foram propostas métricas quantitativas, organizadas em categorias que medem propriedades diferentes das explicações~\cite{hedstrom2023quantus}. As de \emph{fidelidade} verificam se os atributos indicados como importantes são os que mais alteram a saída quando retirados~\cite{bhatt2020}; as de \emph{robustez}, se a explicação muda pouco quando a entrada muda pouco~\cite{yeh2019}; as de \emph{complexidade}, se a importância se concentra em poucos atributos~\cite{bhatt2020}; e as de \emph{randomização}, se a explicação depende do que o modelo aprendeu e da classe explicada~\cite{adebayo2018,sixt2020}.

Em ECG e em séries temporais, essas avaliações mostraram que a qualidade das explicações não é uniforme. A \FC~\cite{bhatt2020}, uma métrica de fidelidade, foi usada para comparar onze métodos de atribuição em quatro bases de ECG~\cite{suh2024}. Quando confrontadas com regiões definidas por diretrizes clínicas, as explicações de métodos que multiplicam o gradiente pela entrada ou subtraem uma referência, entre eles o GradientSHAP, seguiram a amplitude do sinal e deixaram de cobrir regiões de baixa amplitude importantes para o diagnóstico~\cite{gumpfer2026}. Em séries temporais, a eficácia das métricas de perturbação variou muito entre as classes~\cite{baer2025a}, e diferenças simples de amplitude ou de duração do padrão entre as classes bastaram para produzir esse efeito~\cite{baer2025b}.

Esses trabalhos, porém, não relacionam a qualidade das explicações com o desempenho do modelo em cada classe. As comparações de métodos costumam usar classificadores de desempenho alto: o XTSC-Bench avalia apenas modelos com acurácia acima de 90\%~\cite{hollig2023}, e a comparação de Ismail et al.\ relata resultados de modelos com acurácia acima de 95\%~\cite{ismail2020}. Em ECG, um estudo com nove classes usou modelos com AUPRC média acima de 0,85 em oito delas~\cite{arends2026}. Assim, quando o desempenho do modelo é desigual entre as classes, como no CPSC2018, não se sabe se as classes em que o modelo mais acerta são também as de explicações de melhor qualidade.

Essa relação importa para o uso das explicações como apoio ao diagnóstico. Se a qualidade das explicações acompanhar o desempenho, as classes em que o modelo mais erra serão também aquelas em que as explicações menos merecem confiança, e o desempenho poderá servir de indicador da confiabilidade da explicação; se não acompanhar, a qualidade da explicação precisará ser avaliada à parte, classe a classe. A pergunta que orientou este estudo foi, portanto: \emph{a qualidade das explicações GradientSHAP de uma rede neural profunda para classificação de arritmias em ECG acompanha o desempenho (F1) do modelo em cada classe?}

Este estudo comparou, classe a classe, sete métricas de qualidade das explicações com o F1 de uma rede residual treinada no CPSC2018 e mostrou que a resposta depende da propriedade medida. As métricas que envolvem a decisão do modelo entre classes acompanharam o F1 nas classes de melhor e de pior desempenho, enquanto as de fidelidade e de robustez acompanharam o tipo de evidência de cada classe, e não o seu desempenho.
